\documentclass[10pt,a4paper]{amsart}
\usepackage[margin=19mm]{geometry}
\usepackage{amsmath,amssymb,mathtools}
\usepackage[scr=rsfs,cal=euler]{mathalfa}
\usepackage{xfrac}
\usepackage[unicode,hidelinks,bookmarks=false]{hyperref}
\newcommand{\ie}{i.e.\ }
\newcommand{\cf}{cf.\ }
\newcommand{\resp}{respectively\ }
\newcommand{\Subsection}{\subsection}
\providecommand{\nobreakdash}{\nobreak\hbox{-}}
\setlength{\emergencystretch}{2em}
\multlinegap0pt
\usepackage{bm}
\usepackage{bbm}
\usepackage{dsfont}
\usepackage{xspace}
\usepackage{cancel}
\usepackage{mathtools}
\usepackage[shortlabels]{enumitem}
\usepackage{amsthm}
\makeatletter
\@ifundefined{theorem}{\newtheorem{theorem}{Theorem}}{}
\@ifundefined{proposition}{\newtheorem{proposition}[theorem]{Proposition}}{}
\makeatother
\title[Labeled Trees Generating Separable and Locally Finite Ultram]{Labeled Trees Generating Separable and Locally Finite Ultrametrics}
\author{Oleksiy Dovgoshey, Olga Rovenska}
\date{Source: arXiv:2506.03853v2 (CC BY 4.0)}
\begin{document}
\maketitle

\noindent\emph{Source and licence.} Labeled Trees Generating Separable and Locally Finite Ultrametrics. Version arXiv:2506.03853v2 (CC BY 4.0), distributed under the Creative Commons Attribution 4.0 International licence, as recorded in the arXiv licence metadata for this exact version and shown on the abstract page https://arxiv.org/abs/2506.03853v2. Full rights notice: licenses/arxiv-2506.03853-CC-BY-4.0.txt.\par\smallskip

\noindent\textit{Editorial scope.} Counted unit: the Proposition whose source label is \texttt{52dcij} in the article (source0.tex lines 714--726, source label \texttt{52dcij}), delivered together with the article's own complete original proof of it (source0.tex lines 729--819, one \texttt{proof} environment of 91 lines). No mathematics is written, repaired or re-derived: statement and proof are the article's own text line for line. Required context retained uncounted: e11.3 (lines 404--409). Every definition, display and label of those lines is retained unchanged. Omitted from this package and not redistributed: 1--403, 410--713, 727--728, 820--1350. Nothing inside a retained excerpt is omitted or altered: every retained body is delivered exactly as the article's lines are written, so no omission receipt of any kind accompanies this package. Citations: the retained excerpts carry no citation command at all, so no bibliography entry of the article is transcribed and no reference is redistributed. Reference adaptation: none was needed and none was made; every reference inside the delivered unit resolves inside it, so no unresolved reference or citation is produced. Printed slips kept literal: no printed word, symbol, hypothesis or number of the counted statement or of its proof was altered. Layout and class: the article's own class is \texttt{sn-jnl} with packages including graphicx, multirow, amsmath, amssymb, amsfonts, amsthm, mathrsfs, appendix, xcolor, textcomp, manyfoot, booktabs, algorithm, algorithmicx; the wrapper is the lane's pinned \texttt{amsart} editorial wrapper at 10pt on A4, which restates the theorem-like environments the retained text needs and adds the article's own macro definitions that the retained text needs (none), with extra wrapper packages (bm, bbm, dsfont, xspace, cancel, mathtools, enumitem). The delivered numbering is the unit's own. Assets: no retained span contains a figure environment, a graphics-inclusion command, a PDF-page inclusion, a table or a TikZ picture; no image file is copied into this package, the package declares no asset dependency and no image file is redistributed with it. Licence: version arXiv:2506.03853v2 of this article is distributed under the Creative Commons Attribution 4.0 International licence, as recorded in the arXiv licence metadata for this exact version and shown on the abstract page https://arxiv.org/abs/2506.03853v2. A full rights notice is delivered at licenses/arxiv-2506.03853-CC-BY-4.0.txt. Characters: every retained excerpt is pure ASCII, so the delivered engine, XeLaTeX with its default Latin Modern text font, reports no missing character.
\begin{equation}\label{e11.3}
d_l(u, v) = \begin{cases}
0 & \text{if } u = v,\\
\max\limits_{v^{*} \in V(P)} l(v^{*}) & \text{if } u \neq v,
\end{cases}
\end{equation}

\begin{proposition}\label{52dcij}
Let $l_R\colon V(R) \to \mathbb{R}^+$ be a non-degenerate labeling on the vertex set of a ray $R = (v_1, v_2, \ldots, v_n, \ldots)$. Then the following statements are equivalent:

\begin{enumerate}[label=\textit{(\roman*)}, left=0pt]
\item The ultrametric space $(V(R), d_{l_R})$ is locally finite.

\item The limit relation
\begin{equation}\label{tghuaa}
\limsup_{n \to \infty} l_R(v_n) = \infty 
\end{equation}
holds.
\end{enumerate}
\end{proposition}

\begin{proof}

$ (i)$ $\Rightarrow$ $(ii)$. %Let $(V(R), d_{l_R})$ be locally finite.  
Suppose, that the limit superior in \eqref{tghuaa} is a finite %number $k$:
\begin{equation}
\limsup_{n \to \infty} l_R(v_n) = k < \infty.
\end{equation}
Then for each $\varepsilon > 0$, there is $m \in \mathbb{N}$ such that  
\begin{equation*}
l_R(v_n) \leq k + \varepsilon 
\end{equation*}
for all $n \geq m.$ It implies the inequality  
\begin{equation*}
l_R(v_n) \leq \max\{l_R(v_1), \ldots, l_R(v_{m-1}), k + \varepsilon\}
\end{equation*}
for all $n \in \mathbb{N}$. Hence the inequality
\begin{equation}\label{oi45dv}
\sup_{n \in \mathbb{N}} l_R(v_n) < \infty
\end{equation}
holds. Now using \eqref{e11.3} and \eqref{oi45dv}, we obtain
\begin{equation*}
d_{l_R}(u, w) \leq \sup_{n \in \mathbb{N}} l_R(v_n) < \infty
\end{equation*}
for all $u, w \in (V(R),d_{l_R})$. Consequently, $(V(R),d_{l_R})$ is bounded and infinite, contrary to the supposition.

$(ii)$ $\Rightarrow$ $(i)$. %Let $(ii)$ hold. 
Let $A \subseteq V(R)$ be an arbitrary bounded subset.
We must prove that $A$ is finite.
Since $A$ is bounded, there is a constant $c > 0$ such that
\begin{equation}\label{jtfdqw11}
d_{l_R}(u, w) < c 
\end{equation}
for all $u, w \in A$. Using \eqref{tghuaa} we can find an infinite subsequence $(u_{n_k})_{k \in \mathbb{N}}$ of the sequence $(v_n)_{n \in \mathbb{N}}$ such that $n_1 \geq 2$ and
\begin{equation}\label{onhwer12}
l_R(v_{n_k}) \geq c
\end{equation}
for each $k \in \mathbb{N}$. Let us define a sequence $(P_k)_{k \in \mathbb{N}}$ of paths in $R$ by the rule: 
\begin{equation*}
P_1= (v_1, \ldots, v_{n_1})
\end{equation*}
is
the path joining $v_1$ and $v_{n_1}$, and 
\begin{equation*}
P_k = (v_{n_{k-1}}, \ldots, v_{n_k})
\end{equation*}
for every $k \geq 2$.

We claim that a sequence $(P_k)_{k \in \mathbb{N}}$ contains a path $P_{k_0}$ such that
\begin{equation}\label{ettbu531}
A \subseteq V(P_{k_0}). 
\end{equation}

Indeed if the last claim is false, then there exist two distinct paths $P_{k_1}$ and $P_{k_2}$ satisfying the conditions
\begin{equation}\label{ettbu532}
A \cap V(P_{k_1}) \neq \emptyset \neq A \cap V(P_{k_2}), 
\end{equation}
and
\begin{equation}\label{ettbu533}
A \cap V(P_{k_1}) \neq A \cap V(P_{k_2}). 
\end{equation}

Using \eqref{ettbu532} and \eqref{ettbu533}, we can find $a_1, a_2 \in A$ such that
\begin{equation}\label{ettbu534}
a_1 \in V(P_{k_1}) \setminus V(P_{k_2}) \quad \text{and} \quad a_2 \in V(P_{k_2}) \setminus V(P_{k_1}). 
\end{equation}
Without loss of generality, we assume that the inequality
\begin{equation}\label{ettbu535}
k_1 < k_2 
\end{equation}
holds. 
Let $P_{a_1,a_2}$ be a path joining $a_1$ and $a_2$ in $R$. 
Membership relations \eqref{ettbu534} and inequality \eqref{ettbu535} imply 
\begin{equation}\label{ettbu536}
v_{n_{k_1}} \in V(P_{a_1,a_2}). 
\end{equation}
Hence
\begin{equation}\label{ettbu537}
d_{l_R}(a_1, a_2) = \max_{v \in P_{a_1,a_2}} l(v)\geq l(v_{n_{k_1}})
\end{equation}
holds by \eqref{e11.3}, \eqref{ettbu536}.
Now, using \eqref{onhwer12} and \eqref{ettbu537}, we get the inequality $d_{l_R}(a_1, a_2) \geq c$,
which contradicts \eqref{jtfdqw11}.

Thus there is a path $P_{k_0}$ satisfying inclusion \eqref{ettbu531} that implies
\begin{equation*}
|A| \leq |V(P_{k_0})| < \infty.
\end{equation*}

This completes the proof.

\end{proof}

\end{document}
